\section{讨论}

前文描绘了一幅以能力扩张为主线的发展图景，但除能力本身之外，临床证据、评估方法与治理框架三方面仍呈现出明显缺口。

第一类问题集中在人机协作的非线性。Goh 等招募 50 名内科医师在 6 个真实病例上做诊断推理，GPT-4 独立组评分 92\%，但“医生 + GPT-4”组仅 76\%，与传统资源组无显著差异[25]；Bean 等在 1,298 名公众用户的预注册随机对照试验中观察到更剧烈的剪刀差：三款主流 LLM 在脱离用户的封闭设置下识别 10 类病症的准确率为 90.8--99.2\%，但真实公众用户经 LLM 辅助后的病症识别率反而显著低于互联网搜索组[26]。反向证据则来自约束化部署：Pais 等的 MEDIC 将 LLM 严格限定为结构化抽取器，在在线药房前瞻部署中使近失误事件下降 33\%[27]；PreA 在两家三甲医院的随机对照试验中把问诊时长缩短 28.7\%、跨科室协调度提升 113.1\%[20]；心脏专科 AMIE 将显著错误率从 24.3\% 降到 13.1\%[19]。这些证据共同说明，模型独立基准成绩与真实临床效用之间存在系统性的“协作折扣”，而工作流约束、用户素养与共同设计强度，比模型规模本身更能解释这一折扣的方向与大小。

第二类问题在于评估方法学的拓扑滞后。MedHELM 与 29 名临床医生联合构建了 5 大类、22 子类、121 个临床任务及配套 37 个基准，发现所有模型在行政与工作流类别得分最差，恰好对应第一类问题中最易出问题的领域[18]。MetaMedQA 把“我不知道”作为显式选项加入 MedQA 扩展集，结果 12 个主流模型中有 9 个 Unknown Recall 为 0[2]；Kim 等构造的对抗推理题表明，人类医生平均 66\%，最优 LLM 仅 52\%，而医疗专用模型在该任务上甚至接近 0\%[23]。EquityMedQA 与 RadFact 则分别在公平性与多模态接地两个维度上引入新的评估面[24,40]。合并这些结果可以看到，评估科学的当务之急不在于构造更难的题目，而在于扩展评估空间的拓扑，包括任务覆盖、不确定性、公平性、空间接地、被试身份与报告规范等多个维度。

第三类问题潜伏于训练--治理链路的隐性脆弱。Alber 等的数据投毒实验表明，只需在通用 LLM 预训练语料中替换 0.001\% 的训练 token，便足以让 4B 模型在医疗有害内容生成率上显著上升，而所有被投毒模型在 5 项主流医疗基准上的得分与未投毒基线无显著差异[29]。这说明，任何依赖标准基准做上市前评估的监管体系，都可能对这类攻击结构性失效。该脆弱性又与训练数据透明度不足相互放大：系统综述显示 87.7\% 的医疗 LLM 为闭源体系[1]，而 Med-PaLM M、Med-Gemini、MAIRA-2 等代表性模型的训练数据多数私有。监管框架的滞后则被多篇综述和共识反复指出：现行 FDA 与 NMPA 监管基于固定权重，与持续学习、智能体调用、自我更新的 LLM 范式不匹配[3,7,42]。因此，训练数据透明度应当被视为医疗 AI 安全性的必要条件，而不是可选的良好实践。

总体而言，医疗 AI 的能力扩张本身不足以自动保证临床效用、评估有效性与治理合规。患者直面部署还需要超出基准评估的额外维度，例如可读性、健康素养与文化适配[28]。下一阶段的进步很可能不再主要来自更大的模型，而来自把模型与专科底座、外部验证器以及持续审计层组装为可工程化、可监管化的系统。
